%!TEX root = main.tex

Let $ss$ be the statement sequence before applying the function {\tt expand\_cnct} and $ss_{\rm cnct}$ be the sequence of connect statements in $ss$. 

\paragraph*{Specification}
Then the correctness of {\tt expand\_cnct} is specified by the formula 
%
$$\mbox{\tt type\_equivs}(\mbox{\tt expand\_cnct}(ss_{\rm cnct})) \wedge \mbox{\tt all\_ground}(\mbox{\tt expand\_cnct}(ss_{\rm cnct})).$$ 
%
The formula $\mbox{\tt type\_equivs}(\mbox{\tt expand\_cnct}(ss_{\rm cnct}))$ specifies that for each connect statement generated by $\mbox{\tt expand\_cnct}(ss_{\rm cnct})$, the types of its left-hand-side and right-hand-side expressions are equivalent. On the other hand, the formula $\mbox{\tt all\_ground}(\mbox{\tt expand\_cnct}(ss_{\rm cnct}))$ specifies that all the expressions occurring in the connect statements generated by $\mbox{\tt expand\_cnct}(ss_{\rm cnct})$ have ground types. 

More specifically, 
the predicate $\mbox{\tt type\_equivs}(ss')$ is recursively defined and requires that each connect statement  $x <= e$ in  $ss'$ satisfies that $\mbox{\tt type\_equiv}(\mbox{\tt type\_of\_expr}(x), \mbox{\tt type\_of\_expr}(e))$ hods, where ${\tt type\_equiv}(t_1, t_2)$ is the predicate that specifies the type equivalence of $t_1$ and $t_2$. On the other hand, the predicate $\mbox{\tt all\_ground}(ss')$ is recursively defined and requires that each connect statement $x <= e$ in  $ss'$ satisfies that $\mbox{\tt ground}(\mbox{\tt type\_of\_expr}(x)) \wedge \mbox{\tt ground}(\mbox{\tt type\_of\_expr}(e))$ holds, where the predicate $\mbox{\tt ground}(t)$ checks that $t$ is a ground type.   

\paragraph*{Verification}
Recall that {\tt expand\_cnct} expands a connect statement between components of aggregate types by using the function {\tt expand\_expr} to expand each of the left-hand-side and right-hand-side expressions of the statement into a sequence of expressions of ground types. Since $\mbox{\tt expand\_expr}(e)$ expands the expression $e$ by following the structure of $\mbox{\tt type\_of\_expr}(e)$, we define a recursive function $\mbox{\tt expand\_type}(t)$ that expands an aggregate type $t$ into a list of ground types. Then we prove the following two lemmas to prove $ce \vdash \mbox{\tt type\_equivs}(\mbox{\tt expand\_cnct}(ss_{\rm cnct}))$.

\begin{lemma}\label{expand-expr-expand-type}
For every connect statement $x <= e$ in $ss_{\rm cnct}$, 
$$
\begin{array}{l c l}
ce & \vdash & \mbox{\tt type\_equiv}(\mbox{\tt type\_of\_expr}(x), \mbox{\tt type\_of\_expr}(e)) \rightarrow \\
& & \ \ {\tt expand\_type}({\tt type\_of\_expr}(x)) = {\tt expand\_type}({\tt type\_of\_expr}(e)).
\end{array}
$$
\end{lemma}

\begin{lemma}\label{expand-expr-expand-expr-type}
For every expression $e$, 
$$ce \vdash {\tt fold}({\tt type\_of\_expr}, {\tt expand\_expr}(e)) = {\tt expand\_type}\,({\tt type\_of\_expr}(e)).$$
\end{lemma}
Note that in Lemma~\ref{expand-expr-expand-expr-type}, the ${\tt fold}$ function is used to apply the function ${\tt type\_of\_expr}$ to every expression in ${\tt expand\_expr}(e)$ to produce a list of types. 

Finally, $ce \vdash \mbox{\tt all\_ground}(\mbox{\tt expand\_cnct}(ss_{\rm cnct}))$ can be proved by the following lemma. 
\begin{lemma}\label{expand-expr-ground}
For every expression $e$, 
$$ce \vdash {\tt fold}({\tt ground}, {\tt expand\_type}({\tt type\_of\_expr}(e))).$$
\end{lemma}
Note that in Lemma~\ref{expand-expr-ground}, the ${\tt fold}$ function is used to check whether the predicate {\tt ground} is satisfied by every type in ${\tt expand\_type}({\tt type\_of\_expr}(e)))$.


%%%%%% original texts by xiaomu 
%%%%%% original texts by xiaomu 
\hide{
Recall the the ExpandConnects pass is aimed to decompose the connect statements of aggregate
types to ground type connections.
To prove the functional correctness of the ExpandConnects pass, we need to ensure
that the sub elements of the components are connected to their equivalent ground types.
Let $ce^{in}$ be the input component environment,
%% such that $ce^{in}$ is the result of the InferTypes pass,
and $ss^{in}$ be the list of connect statements in the input circuit
diagram $D_M^{in}$.
%% Let $ss^{out}$ be the output list of connect
%% statements after applying the Expand Connections pass.
The input
$D_M^{in}$ is required to be the result of the InferTypes pass, in order
to make sure that the connections are made between equivalent types.

We specify the correctness based on the types of the left-hand side and the right-hand side
expressions in the resulting connect statements.
First, we state that the type equivalence holds for the output statements generated by the ExpandConnects pass
as,
\begin{lemma}\label{expand-conct-type1-ss}
$\forall ce^{in}\; ss^{in},\, ce^{in} \models TypeEquivs\, ss^{in} \Rightarrow TypeEquivs\, [{\tt expand\_cnct}\,s\,|\,s\leftarrow ss^{in}]).$
\end{lemma}
Intuitively, it states that if the type equivalence predicate holds for all the
connect statements in $ss^{in}$,
it holds for the statements produced by applying the {\tt expand\_cnct} \eqref{expand-cnct} function on all the inputs.
More specifically, the type equivalence predicate $TypeEquivs$ holds for a list of statements if
for each individual comparison of the types on the two sides of a connection, the type equivalence
predicate $TypeEquiv$ holds.
If there are $n$ statements $s_1,s_2,\dots ,s_n$ in the sequence $ss$ of statements,
then $TypeEquivs\;ss\,\!=\,\!\bigwedge_{i=0}^{n} TypeEquiv\;s_i$.

The proof can be achieved by induction on the list of statements,
assuming we have the proof for expanding one statement.
Since the function {\tt expand\_cnct} only affect connect statements,
we can state the following lemma:
\begin{lemma}\label{expand-conct-type1-s}
$\forall ce^{in}\;v\;e,\, ce^{in} \models TypeEquiv\,(v<=e) \Rightarrow TypeEquivs\, ({\tt expand\_cnct}\,(v<=e)).$
\end{lemma}
If the predicate $TypeEquiv$ holds for a connect statement $v<=e$ under a component environment,
then the types of $v$ and $e$ are equivalent.
Lemma \ref{expand-conct-type1-s} can be further simplified as,
\begin{lemma}\label{expand-conct-type1-s}
$\forall ce^{in}\;v\;e,\, ce^{in} \models Equiv\,({\tt typeof}\,v,\;{\tt typeof}\,e) \Rightarrow$
$Equivs\,[({\tt typeof}\,l,\,{\tt typeof}\,r)\,|\,l \leftarrow{\tt expand\_expr}\,v,\;r\leftarrow {\tt expand\_expr}\,e]$.
\end{lemma}
If there are $n$ pairs of types $(t_1^l,t_1^r),(t_2^l,t_2^r),\dots ,(t_n^l,t_n^r)$ in the list $ts$,
then $Equivs$ holds for $ts$ if
$\bigwedge_{i=0}^{n} Equiv\;(t_i^l\,t_i^r)$.
In Lemma \ref{expand-conct-type1-s}, the list of pairs are the types obtained by
expanding the reference on the left-hand side
and expanding the expression on the right-hand side, respectively.
If we can build the relation between the aggregate type of an expression and the types of
the corresponding expanded expressions, the above lemma can be proved.
Here we introduce a function {\tt dec\_type} that can decompose the aggregate types
by its structure to get a list of ground types.
A related lemma is,
\begin{lemma}\label{dec-type}
$\forall ce^{in}\, t_1\, t_2,\,ce^{in}\models Equiv\,(t_1,\;t_2) \Rightarrow$
$Equivs\,[(x,\,y)\,|\,x \leftarrow{\tt dec\_type}\,t_1,\;y\leftarrow {\tt dec\_type}\,$ $t_2]$
\end{lemma}
It states that, decomposing a pair of equivalent aggregate types produces two pairwise equivalent type lists.
If we use the {\tt dec\_type} function on the aggregate type of the expression  to be expanded, it
returns the ground types that are identical to the ones produced by {\tt expand\_expr}, as in the following lemma.
\begin{lemma}\label{expand-expr-dec-type}
$\forall ce^{in}\, e,\,ce^{in}\models[{\tt typeof}\,x\,|\,x\leftarrow{\tt expand\_expr}\,e] = {\tt dec\_type}\,({\tt typeof}\,e)$
\end{lemma}
By rewriting Lemma \ref{expand-expr-dec-type} in the goal of Lemma \ref{expand-conct-type1-s}, the proof goal is change to
$Equivs\,[(x,\,y)\,|\,x\leftarrow{\tt dec\_type}\;({\tt typeof}\,v),\;y\leftarrow {\tt dec\_type}\,({\tt typeof}\,e)]$.
Now we can apply Lemma \ref{dec-type} with assigning $t_1$ and $t_2$ with ${\tt typeof}\,v$ and ${\tt typeof}\,e$
to complete the proof.

Secondly, we state that all the expressions after applying the ExpandConnects pass
have ground types.
\begin{lemma}
$\forall ce ^{in}\,ss^{in},\,ce^{in}\models AllGround ({\tt expand\_cnct}\,ss^{in})$
\end{lemma}

The above lemma states the predicate $AllGround$ holds for the generated
statements of {\tt expand\_cnct}.
Similarly, the proof goal will be reduced to prove that the $AllGround$ predicate holds for the expanded expressions.
It is rather easy to use Lemma \ref{expand-expr-dec-type} we had for {\tt dec\_type}.
\fu{To relate the types of the expanded expressions are identical to decompose the type of the original expression.}
By induction on the aggregate type constructors, it is able to prove that decomposed types by
the {\tt dec\_type} function will always return ground types.
}
%%%%%% original texts by xiaomu 
%%%%%% original texts by xiaomu 
